% =====================================================================
%  Appendix D: Coordinate Systems and Formula Reference
% =====================================================================

\chapter{Coordinate Systems and Formula Reference}
\label{app:coordinate-systems-reference}

Coordinate systems are different languages for the same geometry.  Cartesian coordinates
are best for boxes and planes.  Polar and cylindrical coordinates are best for circular
symmetry.  Spherical coordinates are best for balls, shells, cones, and radial fields.

This appendix is a reference.  It is meant to be scanned.

\[
\begin{array}{c|c}
\text{Coordinate system} & \text{Best for} \\ \hline
\text{Cartesian }(x,y,z) & \text{boxes, planes, graphs}\\[1mm]
\text{Polar }(r,\theta) & \text{circles and plane rotation}\\[1mm]
\text{Cylindrical }(r,\theta,z) & \text{cylinders and vertical symmetry}\\[1mm]
\text{Spherical }(\rho,\phi,\theta) & \text{spheres and radial symmetry}
\end{array}
\]

% ---------------------------------------------------------------------
\section{Polar coordinates}
\label{app:D.1}

\subsection*{Conversion formulas}

Polar coordinates describe a point in the plane by radius and angle:
\[
\boxed{
    (r,\theta).
}
\]

Conversion to Cartesian:
\[
\boxed{
    x=r\cos\theta,
    \qquad
    y=r\sin\theta.
}
\]

Conversion to polar:
\[
\boxed{
    r^2=x^2+y^2.
}
\]

\[
\boxed{
    \tan\theta=\frac yx
}
\]
with quadrant care.

\subsection*{Usual ranges}

Usually,
\[
\boxed{
    r\ge0,
    \qquad
    0\le\theta<2\pi.
}
\]

The angle is not unique:
\[
\boxed{
    (r,\theta)=(r,\theta+2\pi k)
}
\]
for every integer \(k\).

If signed radii are allowed, then
\[
\boxed{
    (-r,\theta)=(r,\theta+\pi).
}
\]

\subsection*{Unit vectors}

\[
\boxed{
    \widehat{\mathbf e}_r=(\cos\theta,\sin\theta)
}
\]

\[
\boxed{
    \widehat{\mathbf e}_\theta=(-\sin\theta,\cos\theta)
}
\]

The vector \(\widehat{\mathbf e}_r\) points outward.

The vector \(\widehat{\mathbf e}_\theta\) points in the direction of increasing
\(\theta\).

\subsection*{Unit-vector derivatives}

\[
\boxed{
    \frac{\partial \widehat{\mathbf e}_r}{\partial\theta}
    =
    \widehat{\mathbf e}_\theta
}
\]

\[
\boxed{
    \frac{\partial \widehat{\mathbf e}_\theta}{\partial\theta}
    =
    -\widehat{\mathbf e}_r
}
\]

The polar unit vectors depend on position.

\subsection*{Small displacement}

\[
\boxed{
    d\mathbf r
    =
    \widehat{\mathbf e}_r\,dr
    +
    \widehat{\mathbf e}_\theta\,r\,d\theta.
}
\]

\subsection*{Line element}

\[
\boxed{
    ds^2=dr^2+r^2\,d\theta^2.
}
\]

So
\[
\boxed{
    ds=\sqrt{dr^2+r^2\,d\theta^2}.
}
\]

For a circle \(r=a\):
\[
\boxed{
    ds=a\,d\theta.
}
\]

For a radial line \(\theta=\theta_0\):
\[
\boxed{
    ds=dr.
}
\]

\subsection*{Area element}

Ordinary area:
\[
\boxed{
    dA=r\,dr\,d\theta.
}
\]

Oriented area:
\[
\boxed{
    dx\wedge dy
    =
    r\,dr\wedge d\theta.
}
\]

\subsection*{Common polar curves}

Circle centered at origin:
\[
\boxed{
    r=a.
}
\]

Line through origin:
\[
\boxed{
    \theta=\theta_0.
}
\]

Vertical line:
\[
\boxed{
    x=a
    \quad\Longleftrightarrow\quad
    r\cos\theta=a.
}
\]

Horizontal line:
\[
\boxed{
    y=b
    \quad\Longleftrightarrow\quad
    r\sin\theta=b.
}
\]

Circle through the origin with center on the \(x\)-axis:
\[
\boxed{
    r=2a\cos\theta.
}
\]

Circle through the origin with center on the \(y\)-axis:
\[
\boxed{
    r=2b\sin\theta.
}
\]

\subsection*{Polar integration}

Area of a region:
\[
\boxed{
    \iint_R f(x,y)\,dA
    =
    \iint_{\widetilde R}
    f(r\cos\theta,r\sin\theta)\,r\,dr\,d\theta.
}
\]

Area enclosed by \(r=f(\theta)\), traced once:
\[
\boxed{
    A=\frac12\int_\alpha^\beta f(\theta)^2\,d\theta.
}
\]

\subsection*{Common warnings}

The formula
\[
    \tan\theta=\frac yx
\]
does not determine the quadrant by itself.

The point at the origin has no unique angle.

The factor \(r\) in
\[
    dA=r\,dr\,d\theta
\]
must not be dropped.

Formulas containing \(1/r\) are singular at
\[
    r=0.
\]

% ---------------------------------------------------------------------
\section{Cylindrical coordinates}
\label{app:D.2}

\subsection*{Conversion formulas}

Cylindrical coordinates are polar coordinates in the \(xy\)-plane plus height \(z\).

\[
\boxed{
    (r,\theta,z).
}
\]

Conversion to Cartesian:
\[
\boxed{
    x=r\cos\theta,
    \qquad
    y=r\sin\theta,
    \qquad
    z=z.
}
\]

Conversion from Cartesian:
\[
\boxed{
    r^2=x^2+y^2,
    \qquad
    \tan\theta=\frac yx,
    \qquad
    z=z.
}
\]

\subsection*{Usual ranges}

\[
\boxed{
    r\ge0,
    \qquad
    0\le\theta<2\pi,
    \qquad
    -\infty<z<\infty.
}
\]

\subsection*{Unit vectors}

\[
\boxed{
    \widehat{\mathbf e}_r=(\cos\theta,\sin\theta,0)
}
\]

\[
\boxed{
    \widehat{\mathbf e}_\theta=(-\sin\theta,\cos\theta,0)
}
\]

\[
\boxed{
    \widehat{\mathbf e}_z=(0,0,1).
}
\]

\subsection*{Small displacement}

\[
\boxed{
    d\mathbf r
    =
    \widehat{\mathbf e}_r\,dr
    +
    \widehat{\mathbf e}_\theta\,r\,d\theta
    +
    \widehat{\mathbf e}_z\,dz.
}
\]

\subsection*{Line element}

\[
\boxed{
    ds^2=dr^2+r^2\,d\theta^2+dz^2.
}
\]

\subsection*{Volume element}

Ordinary volume:
\[
\boxed{
    dV=r\,dr\,d\theta\,dz.
}
\]

Oriented volume:
\[
\boxed{
    dx\wedge dy\wedge dz
    =
    r\,dr\wedge d\theta\wedge dz.
}
\]

\subsection*{Coordinate surfaces}

Cylinder:
\[
\boxed{
    r=a.
}
\]

Vertical half-plane:
\[
\boxed{
    \theta=\theta_0.
}
\]

Horizontal plane:
\[
\boxed{
    z=c.
}
\]

\subsection*{Useful equations}

Cylinder of radius \(a\):
\[
\boxed{
    x^2+y^2=a^2
    \quad\Longleftrightarrow\quad
    r=a.
}
\]

Paraboloid:
\[
\boxed{
    z=x^2+y^2
    \quad\Longleftrightarrow\quad
    z=r^2.
}
\]

Cone:
\[
\boxed{
    z^2=x^2+y^2
    \quad\Longleftrightarrow\quad
    z^2=r^2.
}
\]

\subsection*{Cylindrical integration}

\[
\boxed{
    \iiint_E f(x,y,z)\,dV
    =
    \iiint_{\widetilde E}
    f(r\cos\theta,r\sin\theta,z)\,
    r\,dr\,d\theta\,dz.
}
\]

\subsection*{Common warnings}

Cylindrical \(r\) is distance from the \(z\)-axis, not distance from the origin in space.

The factor \(r\) in
\[
    dV=r\,dr\,d\theta\,dz
\]
must not be dropped.

At \(r=0\), the angle \(\theta\) is not unique.

% ---------------------------------------------------------------------
\section{Spherical coordinates}
\label{app:D.3}

\subsection*{Convention used in this book}

Spherical coordinates are
\[
\boxed{
    (\rho,\phi,\theta).
}
\]

Here:

\[
\begin{array}{c|c}
\rho & \text{distance from the origin}\\[1mm]
\phi & \text{angle down from the positive }z\text{-axis}\\[1mm]
\theta & \text{angle around the }z\text{-axis in the }xy\text{-plane}
\end{array}
\]

\subsection*{Conversion formulas}

To Cartesian:
\[
\boxed{
    x=\rho\sin\phi\cos\theta
}
\]

\[
\boxed{
    y=\rho\sin\phi\sin\theta
}
\]

\[
\boxed{
    z=\rho\cos\phi
}
\]

From Cartesian:
\[
\boxed{
    \rho^2=x^2+y^2+z^2.
}
\]

\[
\boxed{
    r=\sqrt{x^2+y^2}=\rho\sin\phi.
}
\]

\[
\boxed{
    z=\rho\cos\phi.
}
\]

\[
\boxed{
    \tan\theta=\frac yx
}
\]
with quadrant care.

\subsection*{Usual ranges}

\[
\boxed{
    \rho\ge0,
    \qquad
    0\le\phi\le\pi,
    \qquad
    0\le\theta<2\pi.
}
\]

\subsection*{Unit vectors}

\[
\boxed{
    \widehat{\mathbf e}_\rho
    =
    (\sin\phi\cos\theta,\sin\phi\sin\theta,\cos\phi)
}
\]

\[
\boxed{
    \widehat{\mathbf e}_\phi
    =
    (\cos\phi\cos\theta,\cos\phi\sin\theta,-\sin\phi)
}
\]

\[
\boxed{
    \widehat{\mathbf e}_\theta
    =
    (-\sin\theta,\cos\theta,0)
}
\]

The vector \(\widehat{\mathbf e}_\rho\) points outward from the origin.

The vector \(\widehat{\mathbf e}_\phi\) points in the direction of increasing \(\phi\).

The vector \(\widehat{\mathbf e}_\theta\) points in the direction of increasing \(\theta\).

\subsection*{Small displacement}

\[
\boxed{
    d\mathbf r
    =
    \widehat{\mathbf e}_\rho\,d\rho
    +
    \widehat{\mathbf e}_\phi\,\rho\,d\phi
    +
    \widehat{\mathbf e}_\theta\,\rho\sin\phi\,d\theta.
}
\]

\subsection*{Line element}

\[
\boxed{
    ds^2
    =
    d\rho^2+\rho^2\,d\phi^2+\rho^2\sin^2\phi\,d\theta^2.
}
\]

\subsection*{Volume element}

Ordinary volume:
\[
\boxed{
    dV=\rho^2\sin\phi\,d\rho\,d\phi\,d\theta.
}
\]

Oriented volume:
\[
\boxed{
    dx\wedge dy\wedge dz
    =
    \rho^2\sin\phi\,
    d\rho\wedge d\phi\wedge d\theta.
}
\]

\subsection*{Coordinate surfaces}

Sphere:
\[
\boxed{
    \rho=a.
}
\]

Cone:
\[
\boxed{
    \phi=\phi_0.
}
\]

Vertical half-plane:
\[
\boxed{
    \theta=\theta_0.
}
\]

\subsection*{Useful equations}

Sphere of radius \(a\):
\[
\boxed{
    x^2+y^2+z^2=a^2
    \quad\Longleftrightarrow\quad
    \rho=a.
}
\]

Upper cone:
\[
\boxed{
    z=\sqrt{x^2+y^2}
    \quad\Longleftrightarrow\quad
    \phi=\frac{\pi}{4}.
}
\]

Plane \(z=c\):
\[
\boxed{
    \rho\cos\phi=c.
}
\]

Cylinder \(x^2+y^2=a^2\):
\[
\boxed{
    \rho\sin\phi=a.
}
\]

\subsection*{Spherical integration}

\[
\boxed{
    \iiint_E f(x,y,z)\,dV
    =
    \iiint_{\widetilde E}
    f(\rho\sin\phi\cos\theta,\rho\sin\phi\sin\theta,\rho\cos\phi)
    \rho^2\sin\phi\,d\rho\,d\phi\,d\theta.
}
\]

\subsection*{Common warnings}

Different books sometimes swap the names \(\theta\) and \(\phi\).  In this book,
\[
    \phi
\]
is measured down from the positive \(z\)-axis.

At \(\rho=0\), the angles are not unique.

At \(\phi=0\) and \(\phi=\pi\), the angle \(\theta\) is not unique.

The factor
\[
    \rho^2\sin\phi
\]
must not be dropped.

% ---------------------------------------------------------------------
\section{Grad, curl, and divergence in common coordinates}
\label{app:D.4}

\subsection*{Cartesian coordinates}

For
\[
    f=f(x,y,z),
\]
\[
\boxed{
    \nabla f
    =
    f_x\,\mathbf i+f_y\,\mathbf j+f_z\,\mathbf k.
}
\]

For
\[
    \mathbf F=(P,Q,R),
\]
\[
\boxed{
    \nabla\cdot\mathbf F
    =
    P_x+Q_y+R_z.
}
\]

\[
\boxed{
    \nabla\times\mathbf F
    =
    \begin{vmatrix}
    \mathbf i&\mathbf j&\mathbf k\\
    \partial_x&\partial_y&\partial_z\\
    P&Q&R
    \end{vmatrix}.
}
\]

That is,
\[
\boxed{
    \nabla\times\mathbf F
    =
    (R_y-Q_z)\mathbf i
    +(P_z-R_x)\mathbf j
    +(Q_x-P_y)\mathbf k.
}
\]

The Laplacian is
\[
\boxed{
    \Delta f
    =
    \nabla\cdot\nabla f
    =
    f_{xx}+f_{yy}+f_{zz}.
}
\]

\subsection*{Plane Cartesian formulas}

For
\[
    \mathbf F=(P,Q),
\]
the scalar curl in the plane is
\[
\boxed{
    \operatorname{curl}_{\mathrm{plane}}\mathbf F
    =
    Q_x-P_y.
}
\]

The associated exterior derivative is
\[
\boxed{
    d(P\,dx+Q\,dy)
    =
    (Q_x-P_y)\,dx\wedge dy.
}
\]

\subsection*{Polar coordinates}

Write
\[
    \mathbf F
    =
    F_r\,\widehat{\mathbf e}_r
    +
    F_\theta\,\widehat{\mathbf e}_\theta.
\]

Gradient:
\[
\boxed{
    \nabla f
    =
    f_r\,\widehat{\mathbf e}_r
    +
    \frac1r f_\theta\,\widehat{\mathbf e}_\theta.
}
\]

Divergence:
\[
\boxed{
    \nabla\cdot\mathbf F
    =
    \frac1r\frac{\partial}{\partial r}(rF_r)
    +
    \frac1r\frac{\partial F_\theta}{\partial\theta}.
}
\]

Scalar curl:
\[
\boxed{
    \operatorname{curl}_{\mathrm{plane}}\mathbf F
    =
    \frac1r\frac{\partial}{\partial r}(rF_\theta)
    -
    \frac1r\frac{\partial F_r}{\partial\theta}.
}
\]

Laplacian:
\[
\boxed{
    \Delta f
    =
    f_{rr}
    +
    \frac1r f_r
    +
    \frac1{r^2}f_{\theta\theta}.
}
\]

\subsection*{Cylindrical coordinates}

Write
\[
    \mathbf F
    =
    F_r\,\widehat{\mathbf e}_r
    +
    F_\theta\,\widehat{\mathbf e}_\theta
    +
    F_z\,\widehat{\mathbf e}_z.
\]

Gradient:
\[
\boxed{
    \nabla f
    =
    f_r\,\widehat{\mathbf e}_r
    +
    \frac1r f_\theta\,\widehat{\mathbf e}_\theta
    +
    f_z\,\widehat{\mathbf e}_z.
}
\]

Divergence:
\[
\boxed{
    \nabla\cdot\mathbf F
    =
    \frac1r\frac{\partial}{\partial r}(rF_r)
    +
    \frac1r\frac{\partial F_\theta}{\partial\theta}
    +
    \frac{\partial F_z}{\partial z}.
}
\]

Curl:
\[
\boxed{
\begin{aligned}
\nabla\times\mathbf F
&=
\frac1r
\left(
    \frac{\partial F_z}{\partial\theta}
    -
    \frac{\partial}{\partial z}(rF_\theta)
\right)\widehat{\mathbf e}_r\\[1mm]
&\quad+
\left(
    \frac{\partial F_r}{\partial z}
    -
    \frac{\partial F_z}{\partial r}
\right)\widehat{\mathbf e}_\theta\\[1mm]
&\quad+
\frac1r
\left(
    \frac{\partial}{\partial r}(rF_\theta)
    -
    \frac{\partial F_r}{\partial\theta}
\right)\widehat{\mathbf e}_z.
\end{aligned}
}
\]

Laplacian:
\[
\boxed{
    \Delta f
    =
    f_{rr}
    +
    \frac1r f_r
    +
    \frac1{r^2}f_{\theta\theta}
    +
    f_{zz}.
}
\]

\subsection*{Spherical coordinates}

Write
\[
    \mathbf F
    =
    F_\rho\,\widehat{\mathbf e}_\rho
    +
    F_\phi\,\widehat{\mathbf e}_\phi
    +
    F_\theta\,\widehat{\mathbf e}_\theta.
\]

Gradient:
\[
\boxed{
    \nabla f
    =
    f_\rho\,\widehat{\mathbf e}_\rho
    +
    \frac1\rho f_\phi\,\widehat{\mathbf e}_\phi
    +
    \frac1{\rho\sin\phi}f_\theta\,\widehat{\mathbf e}_\theta.
}
\]

Divergence:
\[
\boxed{
    \nabla\cdot\mathbf F
    =
    \frac1{\rho^2}
    \frac{\partial}{\partial\rho}(\rho^2F_\rho)
    +
    \frac1{\rho\sin\phi}
    \frac{\partial}{\partial\phi}(\sin\phi\,F_\phi)
    +
    \frac1{\rho\sin\phi}
    \frac{\partial F_\theta}{\partial\theta}.
}
\]

Curl:
\[
\boxed{
\begin{aligned}
\nabla\times\mathbf F
&=
\frac1{\rho\sin\phi}
\left(
    \frac{\partial}{\partial\phi}(\sin\phi\,F_\theta)
    -
    \frac{\partial F_\phi}{\partial\theta}
\right)
\widehat{\mathbf e}_\rho\\[1mm]
&\quad+
\frac1\rho
\left(
    \frac1{\sin\phi}\frac{\partial F_\rho}{\partial\theta}
    -
    \frac{\partial}{\partial\rho}(\rho F_\theta)
\right)
\widehat{\mathbf e}_\phi\\[1mm]
&\quad+
\frac1\rho
\left(
    \frac{\partial}{\partial\rho}(\rho F_\phi)
    -
    \frac{\partial F_\rho}{\partial\phi}
\right)
\widehat{\mathbf e}_\theta.
\end{aligned}
}
\]

Laplacian:
\[
\boxed{
    \Delta f
    =
    f_{\rho\rho}
    +
    \frac2\rho f_\rho
    +
    \frac1{\rho^2}f_{\phi\phi}
    +
    \frac{\cot\phi}{\rho^2}f_\phi
    +
    \frac1{\rho^2\sin^2\phi}f_{\theta\theta}.
}
\]

Equivalently,
\[
\boxed{
    \Delta f
    =
    \frac1{\rho^2}\frac{\partial}{\partial\rho}
    \left(\rho^2 f_\rho\right)
    +
    \frac1{\rho^2\sin\phi}
    \frac{\partial}{\partial\phi}
    \left(\sin\phi\,f_\phi\right)
    +
    \frac1{\rho^2\sin^2\phi}f_{\theta\theta}.
}
\]

\subsection*{Component warning}

In polar, cylindrical, and spherical formulas,
\[
    F_r,\ F_\theta,\ F_\rho,\ F_\phi,\ F_z
\]
are components in the local unit-vector basis, not Cartesian components.

For example,
\[
    \mathbf F=F_r\widehat{\mathbf e}_r+F_\theta\widehat{\mathbf e}_\theta
\]
does not mean
\[
    \mathbf F=(F_r,F_\theta)
\]
in Cartesian coordinates.

\subsection*{Singularity warning}

Polar and cylindrical formulas with \(1/r\) require
\[
    r>0.
\]

Spherical formulas with \(1/\rho\) require
\[
    \rho>0.
\]

Spherical formulas with \(1/\sin\phi\) require
\[
    \sin\phi\ne0.
\]

Coordinate singularities do not always mean the field itself is singular.  They may mean
the coordinate system has broken down.

% ---------------------------------------------------------------------
\section{Area and volume elements}
\label{app:D.5}

\subsection*{Scale factors}

For orthogonal coordinates
\[
    (u_1,u_2,u_3),
\]
the scale factors are
\[
\boxed{
    h_i=\left|\frac{\partial\mathbf r}{\partial u_i}\right|.
}
\]

The line element is
\[
\boxed{
    ds^2=h_1^2\,du_1^2+h_2^2\,du_2^2+h_3^2\,du_3^2.
}
\]

The volume element is
\[
\boxed{
    dV=h_1h_2h_3\,du_1\,du_2\,du_3.
}
\]

\subsection*{Scale factor table}

\[
\begin{array}{c|c|c}
\text{Coordinates} & \text{Scale factors} & \text{Volume or area element} \\ \hline
(r,\theta)\text{ polar} &
h_r=1,\ h_\theta=r &
dA=r\,dr\,d\theta\\[2mm]

(r,\theta,z)\text{ cylindrical} &
h_r=1,\ h_\theta=r,\ h_z=1 &
dV=r\,dr\,d\theta\,dz\\[2mm]

(\rho,\phi,\theta)\text{ spherical} &
h_\rho=1,\ h_\phi=\rho,\ h_\theta=\rho\sin\phi &
dV=\rho^2\sin\phi\,d\rho\,d\phi\,d\theta
\end{array}
\]

\subsection*{Parametric surface}

For a surface
\[
    \mathbf r(u,v),
\]
the vector area element is
\[
\boxed{
    d\mathbf S
    =
    (\mathbf r_u\times\mathbf r_v)\,du\,dv.
}
\]

The scalar surface area element is
\[
\boxed{
    dS=
    |\mathbf r_u\times\mathbf r_v|\,du\,dv.
}
\]

Flux through the oriented surface is
\[
\boxed{
    \iint_S \mathbf F\cdot\widehat{\mathbf n}\,dS
    =
    \iint_D
    \mathbf F(\mathbf r(u,v))
    \cdot
    (\mathbf r_u\times\mathbf r_v)
    \,du\,dv.
}
\]

\subsection*{Graph surface}

For a graph
\[
    z=f(x,y),
\]
use
\[
    \mathbf r(x,y)=(x,y,f(x,y)).
\]

Then
\[
\boxed{
    \mathbf r_x\times\mathbf r_y
    =
    (-f_x,-f_y,1).
}
\]

Upward-oriented vector area:
\[
\boxed{
    d\mathbf S=(-f_x,-f_y,1)\,dx\,dy.
}
\]

Surface area element:
\[
\boxed{
    dS=
    \sqrt{1+f_x^2+f_y^2}\,dx\,dy.
}
\]

\subsection*{Polar area}

\[
\boxed{
    dA=r\,dr\,d\theta.
}
\]

For a polar curve
\[
    r=f(\theta),
\]
area traced once:
\[
\boxed{
    A=\frac12\int_\alpha^\beta f(\theta)^2\,d\theta.
}
\]

\subsection*{Cylindrical surface elements}

For the cylinder
\[
    r=a,
\]
with parameters \((\theta,z)\),
\[
\boxed{
    dS=a\,d\theta\,dz.
}
\]

For the plane
\[
    z=c,
\]
with polar parameters \((r,\theta)\),
\[
\boxed{
    dS=r\,dr\,d\theta.
}
\]

For the half-plane
\[
    \theta=\theta_0,
\]
with parameters \((r,z)\),
\[
\boxed{
    dS=dr\,dz.
}
\]

\subsection*{Spherical surface elements}

For the sphere
\[
    \rho=a,
\]
with parameters \((\phi,\theta)\),
\[
\boxed{
    dS=a^2\sin\phi\,d\phi\,d\theta.
}
\]

For the cone
\[
    \phi=\phi_0,
\]
with parameters \((\rho,\theta)\),
\[
\boxed{
    dS=\rho\sin\phi_0\,d\rho\,d\theta.
}
\]

For the half-plane
\[
    \theta=\theta_0,
\]
with parameters \((\rho,\phi)\),
\[
\boxed{
    dS=\rho\,d\rho\,d\phi.
}
\]

\subsection*{Boundary orientation reminders}

In the plane, positive boundary orientation means:

\[
\boxed{
    \text{walk so the region stays on your left.}
}
\]

For an upward-oriented surface, use the right-hand rule to orient the boundary.

For a solid, the positive boundary orientation is outward.

\subsection*{Forms version of common elements}

Polar:
\[
\boxed{
    dx\wedge dy
    =
    r\,dr\wedge d\theta.
}
\]

Cylindrical:
\[
\boxed{
    dx\wedge dy\wedge dz
    =
    r\,dr\wedge d\theta\wedge dz.
}
\]

Spherical:
\[
\boxed{
    dx\wedge dy\wedge dz
    =
    \rho^2\sin\phi\,
    d\rho\wedge d\phi\wedge d\theta.
}
\]

\subsection*{Common warnings}

Surface area and flux use different objects:

\[
\boxed{
    dS=|\mathbf r_u\times\mathbf r_v|\,du\,dv
}
\]
is unsigned area.

\[
\boxed{
    d\mathbf S=(\mathbf r_u\times\mathbf r_v)\,du\,dv
}
\]
is oriented vector area.

Swapping the parameters \(u\) and \(v\) reverses orientation:
\[
\boxed{
    \mathbf r_v\times\mathbf r_u
    =
    -(\mathbf r_u\times\mathbf r_v).
}
\]

Ordinary area and volume use positive scale factors.

Differential forms keep orientation signs.

% ---------------------------------------------------------------------
\section*{One-page formula summary}

\[
\boxed{
    x=r\cos\theta,\qquad y=r\sin\theta
}
\]

\[
\boxed{
    dA=r\,dr\,d\theta
}
\]

\[
\boxed{
    x=r\cos\theta,\qquad y=r\sin\theta,\qquad z=z
}
\]

\[
\boxed{
    dV=r\,dr\,d\theta\,dz
}
\]

\[
\boxed{
    x=\rho\sin\phi\cos\theta,
    \qquad
    y=\rho\sin\phi\sin\theta,
    \qquad
    z=\rho\cos\phi
}
\]

\[
\boxed{
    dV=\rho^2\sin\phi\,d\rho\,d\phi\,d\theta
}
\]

\[
\boxed{
    \nabla f
    =
    f_r\,\widehat{\mathbf e}_r
    +
    \frac1r f_\theta\,\widehat{\mathbf e}_\theta
    \quad\text{(polar)}
}
\]

\[
\boxed{
    \nabla f
    =
    f_r\,\widehat{\mathbf e}_r
    +
    \frac1r f_\theta\,\widehat{\mathbf e}_\theta
    +
    f_z\,\widehat{\mathbf e}_z
    \quad\text{(cylindrical)}
}
\]

\[
\boxed{
    \nabla f
    =
    f_\rho\,\widehat{\mathbf e}_\rho
    +
    \frac1\rho f_\phi\,\widehat{\mathbf e}_\phi
    +
    \frac1{\rho\sin\phi}f_\theta\,\widehat{\mathbf e}_\theta
    \quad\text{(spherical)}
}
\]

\[
\boxed{
    \nabla\cdot\mathbf F
    =
    \frac1r\frac{\partial}{\partial r}(rF_r)
    +
    \frac1r\frac{\partial F_\theta}{\partial\theta}
    +
    \frac{\partial F_z}{\partial z}
    \quad\text{(cylindrical)}
}
\]

\[
\boxed{
    \nabla\cdot\mathbf F
    =
    \frac1{\rho^2}
    \frac{\partial}{\partial\rho}(\rho^2F_\rho)
    +
    \frac1{\rho\sin\phi}
    \frac{\partial}{\partial\phi}(\sin\phi\,F_\phi)
    +
    \frac1{\rho\sin\phi}
    \frac{\partial F_\theta}{\partial\theta}
    \quad\text{(spherical)}
}
\]

\[
\boxed{
    dS=|\mathbf r_u\times\mathbf r_v|\,du\,dv
}
\]

\[
\boxed{
    d\mathbf S=(\mathbf r_u\times\mathbf r_v)\,du\,dv
}
\]

\[
\boxed{
    dS=\sqrt{1+f_x^2+f_y^2}\,dx\,dy
    \quad\text{for }z=f(x,y)
}
\]

\[
\boxed{
    dx\wedge dy=r\,dr\wedge d\theta
}
\]

\[
\boxed{
    dx\wedge dy\wedge dz
    =
    \rho^2\sin\phi\,
    d\rho\wedge d\phi\wedge d\theta
}
\]